\documentclass[11pt,a4paper]{article}

\usepackage[margin=1in]{geometry}
\usepackage{hyperref}
\usepackage{listings}
\usepackage{xcolor}
\usepackage{booktabs}
\usepackage{tabularx}
\usepackage{amsmath}
\usepackage{amssymb}
\usepackage{titlesec}
\usepackage{enumitem}
\usepackage{fancyhdr}
\usepackage{graphicx}
\usepackage{microtype}

% ── Colors ──────────────────────────────────────────────────────────────────
\definecolor{codegreen}{rgb}{0.1,0.5,0.1}
\definecolor{codegray}{rgb}{0.4,0.4,0.4}
\definecolor{codepurple}{rgb}{0.5,0,0.5}
\definecolor{codeblue}{rgb}{0,0.3,0.7}
\definecolor{backcolour}{rgb}{0.97,0.97,0.97}
\definecolor{addgreen}{rgb}{0.85,1,0.85}
\definecolor{remred}{rgb}{1,0.85,0.85}
\definecolor{accentblue}{rgb}{0.12,0.33,0.60}

% ── Code listing style ───────────────────────────────────────────────────────
\lstdefinestyle{pycode}{
    backgroundcolor=\color{backcolour},
    commentstyle=\color{codegreen}\itshape,
    keywordstyle=\color{codeblue}\bfseries,
    numberstyle=\tiny\color{codegray},
    stringstyle=\color{codepurple},
    basicstyle=\ttfamily\footnotesize,
    breakatwhitespace=false,
    breaklines=true,
    captionpos=b,
    keepspaces=true,
    numbers=left,
    numbersep=6pt,
    showspaces=false,
    showstringspaces=false,
    showtabs=false,
    tabsize=4,
    language=Python,
    frame=single,
    rulecolor=\color{gray!40},
    xleftmargin=12pt,
    xrightmargin=4pt,
}
\lstset{style=pycode}

% ── Header / Footer ──────────────────────────────────────────────────────────
\pagestyle{fancy}
\fancyhf{}
\fancyhead[L]{\small\textcolor{accentblue}{\textbf{RL Market Maker — Code Changes Report}}}
\fancyhead[R]{\small\textcolor{gray}{CSCI 566}}
\fancyfoot[C]{\thepage}
\renewcommand{\headrulewidth}{0.4pt}

% ── Section formatting ───────────────────────────────────────────────────────
\titleformat{\section}{\large\bfseries\color{accentblue}}{\thesection.}{0.5em}{}[\titlerule]
\titleformat{\subsection}{\normalsize\bfseries\color{accentblue!80}}{\thesubsection}{0.5em}{}

% ── Diff-style macros ────────────────────────────────────────────────────────
\newcommand{\added}[1]{\colorbox{addgreen}{\texttt{\small #1}}}
\newcommand{\removed}[1]{\colorbox{remred}{\texttt{\small #1}}}
\newcommand{\filename}[1]{\texttt{\textbf{\color{accentblue}#1}}}

% ─────────────────────────────────────────────────────────────────────────────
\begin{document}

\begin{titlepage}
    \centering
    \vspace*{2cm}
    {\Huge\bfseries\color{accentblue} RL Market Maker\\[0.4em]}
    {\LARGE Code Changes Report\\[1em]}
    \rule{\linewidth}{1pt}\\[1em]
    {\large\textit{Maker-Taker Fee Integration, PPO Upgrade,\\
    and Bilateral Market-Making Fixes}}\\[2em]
    \rule{\linewidth}{0.5pt}\\[2em]
    {\large CSCI 566 — Final Project\\[0.5em]}
    {\normalsize University of Southern California\\[2em]}
    {\normalsize\textcolor{gray}{April 2026}}
    \vfill
    \begin{abstract}
        \noindent
        This report documents six categories of code changes applied to the
        \texttt{rl\_marketmaker} repository.  The primary goal is to introduce
        a realistic \textbf{maker-taker fee structure} that incentivises
        the reinforcement-learning agent to prefer passive limit orders over
        aggressive market orders, and to upgrade the PPO training loop to
        production-grade standards.  Supporting changes fix bilateral
        market-making order-tracking bugs, remove Colab-specific boilerplate,
        and relax assertions that were incompatible with short-selling.
    \end{abstract}
\end{titlepage}

\tableofcontents
\newpage

% ─────────────────────────────────────────────────────────────────────────────
\section{Overview of Changed Files}

\begin{table}[h!]
\centering
\renewcommand{\arraystretch}{1.3}
\begin{tabularx}{\linewidth}{lX}
\toprule
\textbf{File} & \textbf{Summary of Changes} \\
\midrule
\filename{config/config.py} & Added \texttt{fee\_config} with maker rebate and taker fee values \\
\filename{config/\_\_init\_\_.py} & Exported \texttt{fee\_config} from the config package \\
\filename{simulation/market\_gym.py} & Integrated fee config into \texttt{Market} environment; propagates fees to execution agent \\
\filename{simulation/agents.py} & Added fee accounting to \texttt{ExecutionAgent}; relaxed short-selling assertions \\
\filename{limit\_order\_book/limit\_order\_book.py} & Fixed stale order-ID tracking on cancellation for bilateral MM \\
\filename{rl\_files/actor\_critic.py} & Upgraded PPO to standard clipped loss, gradient clipping, KL tracking; restored production hyperparameters \\
\filename{bilateral\_mm\_agent.ipynb} & Ported notebook from Google Colab to local execution \\
\bottomrule
\end{tabularx}
\caption{Files modified in this changeset}
\end{table}

% ─────────────────────────────────────────────────────────────────────────────
\section{\texttt{config/config.py} — Maker-Taker Fee Configuration}

\subsection{Motivation}

Real exchanges charge takers (agents submitting market orders that
immediately consume liquidity) and reward makers (agents posting limit
orders that add liquidity).  Without this fee structure the RL agent
has no incentive to behave as a true \emph{market maker}; it may prefer
the speed and certainty of market orders.

\subsection{What Was Added}

A new \texttt{fee\_config} dictionary was appended at the end of the file:

\begin{lstlisting}
fee_config = {}
fee_config['maker_rebate'] = 0.2  # reward units per filled lot / initial_volume
fee_config['taker_fee']    = 0.3  # reward units per filled lot / initial_volume
\end{lstlisting}

\subsection{Calibration Rationale}

The fees are expressed in \emph{reward space} (not ticks), so they scale
consistently regardless of price level.

\begin{itemize}[leftmargin=*]
    \item Reference price = 1000, initial volume = 20 lots.
    \item 1 lot filled at 1 tick profit $\Rightarrow$ reward contribution
          $= \frac{1}{20} \times 1 = 0.05$ reward units.
    \item \textbf{Taker fee} $= 0.3$: 1 lot costs $\frac{0.3}{20} = 0.015$
          (30\% of a 1-tick profit — realistic exchange fee).
    \item \textbf{Maker rebate} $= 0.2$: 1 lot earns $\frac{0.2}{20} = 0.010$
          (20\% of a 1-tick profit — typical rebate tier).
    \item Net exchange revenue per lot $= \frac{0.3 - 0.2}{20} = 0.005$
          (approximately 1 basis-point spread, consistent with real markets).
\end{itemize}

Over a full 20-lot episode this translates to:
\[
    \Delta\text{reward}_{\text{all market orders}} = -0.3
    \qquad
    \Delta\text{reward}_{\text{all limit fills}} = +0.2
\]
A 0.5 reward-unit swing provides a strong but not overwhelming incentive
to be the passive side of every trade.  Setting both values to \texttt{0.0}
disables fees completely and preserves backward compatibility.

% ─────────────────────────────────────────────────────────────────────────────
\section{\texttt{config/\_\_init\_\_.py} — Package Export}

A single line was added to expose \texttt{fee\_config} outside the package:

\begin{lstlisting}
from .config import (
    ...
    fee_config,   # <-- added
)

__all__ = [
    ...
    "fee_config", # <-- added
]
\end{lstlisting}

This allows any module that does \lstinline{from config import fee_config} to
access the defaults without importing the full \texttt{config.py} file directly.

% ─────────────────────────────────────────────────────────────────────────────
\section{\texttt{simulation/market\_gym.py} — Market Environment Integration}

\subsection{Importing Fee Defaults}

\begin{lstlisting}
from config import ..., fee_config  # fee_config added to import
\end{lstlisting}

\subsection{Initialisation (\texttt{\_\_init\_\_})}

The \texttt{Market} constructor now reads fee parameters from the external
\texttt{config} dict passed at construction time, falling back to the
global \texttt{fee\_config} defaults if not provided:

\begin{lstlisting}
self.maker_rebate = float(config.get('maker_rebate', fee_config['maker_rebate']))
self.taker_fee    = float(config.get('taker_fee',    fee_config['taker_fee']))
self.cumulative_fees_paid     = 0.0
self.cumulative_rebates_earned = 0.0
\end{lstlisting}

This design means the training script can override fees per-environment
via the config dict without touching global defaults.

\subsection{Agent Injection (\texttt{reset})}

When the market resets and attaches the execution agent, it calls:

\begin{lstlisting}
agent.set_fees(self.maker_rebate, self.taker_fee)
\end{lstlisting}

so the agent always has the correct fee rates before an episode begins.
The market also resets its own fee accumulators:

\begin{lstlisting}
self.cumulative_fees_paid      = 0.0
self.cumulative_rebates_earned = 0.0
\end{lstlisting}

\subsection{Removed Assertion}

\begin{lstlisting}
# Before:
if terminated and not self.circuit_breaker_triggered:
    assert self.agents[self.execution_agent_id].volume == 0

# After: assertion removed
\end{lstlisting}

The bilateral market-making agent can hold a non-zero (or even negative)
position at episode end, so forcing \texttt{volume == 0} was incorrect.

\subsection{Info Dictionary}

Three new keys are reported in the \texttt{info} dict at each step:

\begin{lstlisting}
'cumulative_fees_paid':      float(...),
'cumulative_rebates_earned': float(...),
'net_fee_impact':            float(rebates - fees),
\end{lstlisting}

These allow TensorBoard or downstream analysis to track the fee regime's
real impact on agent profitability.

% ─────────────────────────────────────────────────────────────────────────────
\section{\texttt{simulation/agents.py} — ExecutionAgent Fee Accounting}

\subsection{New \texttt{set\_fees} Method}

\begin{lstlisting}
def set_fees(self, maker_rebate: float, taker_fee: float):
    """Set maker-taker fee structure. Called by Market environment."""
    self.maker_rebate = float(maker_rebate)
    self.taker_fee    = float(taker_fee)
\end{lstlisting}

Fee values default to \texttt{0.0} in \texttt{\_\_init\_\_}, so existing
code that never calls \texttt{set\_fees} is unaffected.

\subsection{Taker Fee (Market Orders)}

When the agent's market order is filled:

\begin{lstlisting}
filled_vol = int(fill_message.filled_volume)
taker_fee_cost = self.taker_fee * filled_vol / self.initial_volume

# Buy side (volume increases):
pnl    = self.get_reward(fill_message.execution_price, filled_vol)
reward -= pnl              # inventory cost
reward -= taker_fee_cost   # exchange fee

# Sell side (volume decreases):
pnl    = self.get_reward(fill_message.execution_price, filled_vol)
reward += pnl              # inventory relief
reward -= taker_fee_cost   # exchange fee (both sides pay)

self.cumulative_fees_paid += taker_fee_cost
\end{lstlisting}

Note that taker fees apply symmetrically to buy \emph{and} sell market
orders: the agent pays every time it demands liquidity.

\subsection{Maker Rebate (Limit Orders)}

When the agent's resting limit order is passively filled:

\begin{lstlisting}
maker_rebate_reward = self.maker_rebate * volume / self.initial_volume

# The rebate is added to reward on top of the PnL term
reward += pnl
reward += maker_rebate_reward
self.cumulative_rebates_earned += maker_rebate_reward
\end{lstlisting}

\subsection{Relaxed Assertions for Bilateral MM}

Several assertions that assumed the agent could only hold a non-negative
inventory were removed or replaced:

\begin{itemize}[leftmargin=*]
    \item \textbf{Removed}: \lstinline{assert self.volume >= 0}
          — the bilateral agent may go short.
    \item \textbf{Removed}: \lstinline{assert self.volume >= 0, f"volume ... < 0"}
          — same reason (duplicate assert with message).
    \item \textbf{Removed}: \lstinline{assert self.active_volume <= self.volume}
          — relaxed; order tracking counters are for bookkeeping only.
\end{itemize}

% ─────────────────────────────────────────────────────────────────────────────
\section{\texttt{limit\_order\_book/limit\_order\_book.py} — Cancellation Bug Fix}

\subsection{Problem}

When an order was cancelled, the LOB correctly removed it from
\texttt{price\_map} and \texttt{order\_map}, but the bilateral
market-making data structures \texttt{agent\_bid\_orders} and
\texttt{agent\_ask\_orders} (sets of active order IDs per agent) were
\emph{not} updated.  This left stale IDs in the sets, causing the MM
agent to believe it had active quotes when it did not.

\subsection{Fix}

Four lines were added after the existing removal logic:

\begin{lstlisting}
self.order_map_by_agent[order.agent_id].remove(id)

# BILATERAL MM: Deregister cancelled order from bid/ask tracking
if order.agent_id in self.agent_bid_orders:
    self.agent_bid_orders[order.agent_id].discard(id)
    self.agent_ask_orders[order.agent_id].discard(id)
\end{lstlisting}

\texttt{discard} (rather than \texttt{remove}) is used intentionally so
that no \texttt{KeyError} is raised if the ID is not present in one of
the two sets (e.g.\ the order was on the opposite side).

% ─────────────────────────────────────────────────────────────────────────────
\section{\texttt{rl\_files/actor\_critic.py} — PPO Training Upgrade}

\subsection{Hyperparameter Restoration}

During debugging, many hyperparameters were reduced to toy values.
These have been restored to production settings:

\begin{table}[h!]
\centering
\renewcommand{\arraystretch}{1.3}
\begin{tabular}{lccc}
\toprule
\textbf{Parameter} & \textbf{Debug Value} & \textbf{Production Value} & \textbf{Reason} \\
\midrule
\texttt{num\_envs}       & 1    & 128  & Parallel sample collection \\
\texttt{num\_steps}      & 10   & 100  & Longer rollouts, better credit assignment \\
\texttt{num\_minibatches}& 1    & 4    & Standard PPO mini-batch update \\
\texttt{update\_epochs}  & 1    & 4    & Multiple passes per rollout \\
\texttt{clip\_coef}      & 0.5  & 0.2  & Standard PPO clipping (Schulman et al.) \\
\texttt{ent\_coef}       & 0.0  & 0.01 & Light entropy bonus for exploration \\
\texttt{total\_timesteps}& 20   & $200\times128\times100$ & Full training run \\
\bottomrule
\end{tabular}
\caption{Hyperparameter changes in \texttt{actor\_critic.py}}
\end{table}

\subsection{New Hyperparameters}

\begin{lstlisting}
max_grad_norm: float = 0.5   # gradient clipping threshold
target_kl:     float = None  # early-stop epoch if KL exceeds threshold
maker_rebate:  float = 0.2   # fee passthrough from Args to env config
taker_fee:     float = 0.3
\end{lstlisting}

\subsection{Standard PPO Clipped Surrogate Loss}

The original code used a simplified policy gradient loss without the
clipping ratio.  This was replaced with the canonical PPO objective
from Schulman et al.\ (2017):

\begin{lstlisting}
# Before (non-standard):
pg_loss = -mb_advantages * newlogprob
pg_loss = pg_loss.mean()

# After (standard PPO):
ratio    = (newlogprob - b_logprobs[mb_inds]).exp()
pg_loss1 = -mb_advantages * ratio
pg_loss2 = -mb_advantages * torch.clamp(ratio,
                                         1 - args.clip_coef,
                                         1 + args.clip_coef)
pg_loss  = torch.max(pg_loss1, pg_loss2).mean()
\end{lstlisting}

The clipping prevents excessively large policy updates, which was a
key stability problem during long training runs.

\subsection{Gradient Clipping}

\begin{lstlisting}
nn.utils.clip_grad_norm_(agent.parameters(), args.max_grad_norm)
\end{lstlisting}

This bounds the gradient norm before the optimiser step, preventing
the policy from collapsing during periods of high reward variance
(common in market-making tasks).

\subsection{KL Divergence Tracking and Early Stop}

The approximate KL divergence between the old and new policy is now
computed and logged after each mini-batch:

\begin{lstlisting}
old_approx_kl = (-logratio).mean()
approx_kl     = ((ratio - 1) - logratio).mean()
clipfracs    += [((ratio - 1.0).abs() > args.clip_coef).float().mean().item()]
\end{lstlisting}

If \texttt{target\_kl} is set, the inner epoch loop breaks early when
the KL exceeds the threshold, preventing policy over-fitting within a
single rollout batch:

\begin{lstlisting}
if args.target_kl is not None and approx_kl > args.target_kl:
    break
\end{lstlisting}

All three diagnostics are written to TensorBoard:

\begin{lstlisting}
writer.add_scalar("losses/old_approx_kl", old_approx_kl.item(), global_step)
writer.add_scalar("losses/approx_kl",     approx_kl.item(),     global_step)
writer.add_scalar("losses/clipfrac",      np.mean(clipfracs),   global_step)
\end{lstlisting}

\subsection{\texttt{AsyncVectorEnv} $\to$ \texttt{SyncVectorEnv}}

\begin{lstlisting}
# Before:
envs = gym.vector.AsyncVectorEnv(env_fns=env_fns)

# After:
envs = gym.vector.SyncVectorEnv(env_fns=env_fns)
\end{lstlisting}

\texttt{AsyncVectorEnv} spawns separate processes and uses a different
\texttt{info} dict format (\texttt{final\_info} list) that was causing
episode-return collection to silently drop data.
\texttt{SyncVectorEnv} runs environments sequentially in the same
process, uses the simpler \texttt{\_key} mask convention, and is easier
to debug.

\subsection{Fee Config Passthrough}

The environment config dicts now carry fee parameters so every parallel
environment receives the same fee regime:

\begin{lstlisting}
base_cfg = {
    'market_env':   args.env_type,
    'execution_agent': 'rl_agent',
    'volume':       args.num_lots,
    'terminal_time': args.terminal_time,
    'time_delta':   args.time_delta,
    'drop_feature': args.drop_feature,
    'maker_rebate': args.maker_rebate,   # <-- new
    'taker_fee':    args.taker_fee,      # <-- new
}
configs = [{**base_cfg, 'seed': args.seed + s} for s in range(args.num_envs)]
\end{lstlisting}

\subsection{\texttt{get\_value} Helper on \texttt{AgentLogisticNormal}}

\begin{lstlisting}
def get_value(self, x):
    return self.critic(x)
\end{lstlisting}

A small convenience method was added so the value head can be called
without unpacking the full actor output.  Used in the GAE advantage
computation bootstrapping step.

% ─────────────────────────────────────────────────────────────────────────────
\section{\texttt{bilateral\_mm\_agent.ipynb} — Local Execution Port}

The notebook was originally designed for Google Colab (with Google Drive
mounting, \texttt{pip install}, and hard-coded \texttt{/content/} paths).
It has been ported to run on a local machine:

\begin{table}[h!]
\centering
\renewcommand{\arraystretch}{1.3}
\begin{tabularx}{\linewidth}{lX}
\toprule
\textbf{Colab behaviour} & \textbf{Local behaviour} \\
\midrule
Mount Google Drive at \texttt{/content/drive} & Uses \texttt{./checkpoints/} in the working directory \\
\texttt{pip install torch ...} & Verifies that dependencies are already installed \\
\texttt{git clone} into \texttt{/content/} & Detects the repo via \texttt{os.getcwd()} \\
\texttt{git pull} on startup & Skipped (local repo managed manually) \\
Hard-coded \texttt{/content/rtle\_parallelized} path & \texttt{repo\_dir = os.getcwd()} \\
\bottomrule
\end{tabularx}
\caption{Notebook environment changes}
\end{table}

All cells that previously raised \texttt{ModuleNotFoundError} (because
\texttt{sortedcontainers} was missing in the Colab kernel before the
install cell ran) now execute cleanly, and outputs have been updated to
reflect the successful local run.

% ─────────────────────────────────────────────────────────────────────────────
\section{Summary}

\begin{table}[h!]
\centering
\renewcommand{\arraystretch}{1.4}
\begin{tabularx}{\linewidth}{lX}
\toprule
\textbf{Change} & \textbf{Impact} \\
\midrule
Maker-taker fee config & Agent now has a measurable incentive to post
                         limit orders rather than market orders.
                         The $\pm 0.5$ reward swing over an episode is
                         significant enough to shape policy without
                         dominating the inventory PnL signal. \\
\addlinespace
Fee integration in Market + Agent & Fees are applied automatically in
                                    every episode with zero code changes
                                    needed in the RL loop. \\
\addlinespace
LOB cancellation bug fix & Eliminates phantom active-quote counts that
                            would cause the bilateral MM agent to
                            miscalculate its spread exposure. \\
\addlinespace
PPO clipped loss + grad clip & Resolves training instability from large
                               policy updates; aligns implementation with
                               the Schulman et al.\ (2017) reference. \\
\addlinespace
Production hyperparameters & Restores the learning signal from 128
                             parallel environments running full 100-step
                             rollouts instead of toy single-env runs. \\
\addlinespace
SyncVectorEnv & Fixes silent episode-return data loss caused by the
                \texttt{AsyncVectorEnv} info dict format mismatch. \\
\addlinespace
Notebook local port & Enables reproducible local experiments without
                      Colab dependency. \\
\bottomrule
\end{tabularx}
\caption{Summary of changes and their impact}
\end{table}

\vfill
\noindent\textcolor{gray}{\small
Generated for CSCI 566 Final Project — RL Market Maker.
Compile with \texttt{pdflatex changes\_report.tex} or upload directly to Overleaf.}

\end{document}
